% Exemplo: química com mhchem (fórmulas e equações), chemfig (fórmula estrutural da cafeína) e
% siunitx (grandezas com unidade: massa, volume, concentração e massa molar).
% Compilar: latexmk -pdf quimica.tex              Testado com: mhchem 4, chemfig 1.6, TeX Live 2024 ou mais recente
%%BEGIN:quimicapreambulo
\documentclass[12pt, a4paper]{article}
\usepackage[T1]{fontenc}
\usepackage[brazilian]{babel}
\usepackage[version=4]{mhchem}   % \ce{...}: fórmulas e equações químicas
\usepackage{chemfig}             % \chemfig{...}: fórmulas estruturais
\usepackage{siunitx}             % \qty{valor}{unidade}: grandezas com unidade
\sisetup{output-decimal-marker = {,}, per-mode = symbol}   % vírgula decimal e mol/L
%%END:quimicapreambulo

\begin{document}
\pagestyle{empty}

%%BEGIN:quimicaequacoes
Combustão completa do metano:
\begin{center}
  \ce{CH4 + 2 O2 -> CO2 + 2 H2O}
\end{center}
Equilíbrio químico: \ce{N2 + 3 H2 <=> 2 NH3}.
Íons e hidratos: \ce{SO4^2-}, \ce{Fe^3+} e \ce{CuSO4 * 5 H2O}.
%%END:quimicaequacoes

%%BEGIN:quimicaestrutura
A cafeína, \ce{C8H10N4O2}, substância indispensável à realização deste livro:
\begin{center}
  \chemfig{*6((=O)-N(-CH_3)-*5(-N=-N(-CH_3)-=)--(=O)-N(-H_3C)-)}
\end{center}
%%END:quimicaestrutura

%%BEGIN:quimicagrandezas
Um café expresso de \qty{30}{\milli\litre} tem cerca de \qty{60}{\milli\gram} de cafeína
(massa molar \qty{194,19}{\gram\per\mole}), ou seja, \qty{2}{\gram\per\litre}, cerca de
\qty{0,01}{\mole\per\litre}.
%%END:quimicagrandezas

\end{document}
